\documentclass{article}
\usepackage[T1]{fontenc}
\usepackage{lmodern}
\usepackage{graphicx}
\usepackage{amsmath,amssymb}
\usepackage[a4paper,margin=2cm]{geometry}
\usepackage[absolute,overlay]{textpos}
\setlength{\TPHorizModule}{1mm}
\setlength{\TPVertModule}{1mm}
\usepackage[indonesian]{babel}
\usepackage{hyperref}
\setlength{\emergencystretch}{2em}
% unit: Halaman 72

\begin{document}

% === Halaman 72 (python/native) ===

Ukuran blok dan kunci

• Ukuran blok pesan (n) yang diproses selama enkripsi/dekripsi selalu 

tetap.

   Contoh: DES (n = 64), AES (n = 128, 192, 256)

• Makin besar ukuran blok mengarah ke efek diffusion yang lebih besar.

• Ukuran kunci (m) bisa sama dengan ukuran blok atau berbeda dengan 

ukuran blok

   Contoh: DES (m = 56), AES (m = 128, 192, 256)

• Makin besar ukuran kunci mengarah ke efek confusion yang lebih 

besar dan lebih tahan terhadap serangan brute force.

72

\end{document}
